\section{Requirements and design constraints}
\label{sec:requirements}

This section states the requirements the system meets and the constraints that shaped it, so the
chapters that follow read as the response to them. The final column of each table points to the
section where the requirement is addressed.

\subsection{Functional requirements}

\Cref{tab:fr} lists what the system must do for an owner.

\begin{table}[ht]
  \centering
  \caption{Functional requirements.}
  \label{tab:fr}
  \begin{tabularx}{\textwidth}{@{}l X l@{}}
    \toprule
    \textbf{ID} & \textbf{Requirement} & \textbf{Addressed in} \\
    \midrule
    FR-1 & Recognise a cat's everyday behaviours (eating, drinking, activity, rest, grooming) from a
           collar-worn sensor. & \cref{sec:on-device-ai} \\
    FR-2 & Perform classification on the wearable itself, in real time, without streaming raw sensor
           data off the cat. & \cref{sec:on-device-ai} \\
    FR-3 & Let the owner define the behaviour set; do not hard-code it in firmware. &
           \cref{sec:training-pipeline} \\
    FR-4 & Provide a closed training loop (capture, label, train, deploy) driven entirely from the
           dashboard, with no local toolchain. & \cref{sec:training-pipeline} \\
    FR-5 & Deliver trained models to the collar over the air, with no firmware reflash. &
           \cref{sec:firmware-station,sec:protocols} \\
    FR-6 & Present the owner a live view and a longitudinal history with health-oriented analytics. &
           \cref{sec:dashboard} \\
    FR-7 & Support multiple owners, each seeing only their own cats and devices. &
           \cref{sec:cloud-backend} \\
    FR-8 & Offer a public, read-only way to explore the product without an account. &
           \cref{sec:dashboard} \\
    \bottomrule
  \end{tabularx}
\end{table}

\subsection{Non-functional requirements}

\Cref{tab:nfr} lists the quality attributes and limits the design must meet, the constraints behind
many of the engineering choices that follow.

\begin{table}[ht]
  \centering
  \caption{Non-functional requirements and quality attributes.}
  \label{tab:nfr}
  \begin{tabularx}{\textwidth}{@{}l X l@{}}
    \toprule
    \textbf{ID} & \textbf{Requirement} & \textbf{Addressed in} \\
    \midrule
    NFR-1 & \textbf{Power.} The collar must run for a useful time on a collar-sized cell, which
            precludes Wi-Fi on the wearable and bounds the duty cycle of the audio path. &
            \cref{sec:hardware,sec:on-device-ai} \\
    NFR-2 & \textbf{Footprint.} Inference must fit the microcontroller's RAM alongside the BLE stack. &
            \cref{sec:on-device-ai} \\
    NFR-3 & \textbf{Privacy.} Raw audio must never leave the collar in normal operation; data must be
            scoped per owner. & \cref{sec:security-privacy} \\
    NFR-4 & \textbf{Credential safety.} Battery devices must not hold the owner's password; device
            authorisation must be revocable. & \cref{sec:security-privacy} \\
    NFR-5 & \textbf{Reproducibility.} A third party must be able to build and run the system from the
            repository without contacting the authors. & \cref{sec:cloud-backend} \\
    NFR-6 & \textbf{Openness.} Hardware designs and code are released under open licences. &
            \cref{sec:hardware} \\
    \bottomrule
  \end{tabularx}
\end{table}

\subsection{Constraints}

\begin{description}[leftmargin=1.6em,style=nextline]
  \item[Radio.] The collar is Bluetooth Low Energy only; BLE is low-power but short-range and is not
        an internet transport, so a mains-powered station is required as the gateway
        (\cref{sec:architecture}).
  \item[Compute.] On-device inference runs on a Cortex-M4F with a few hundred kilobytes of RAM and no
        operating-system heap budget to spare, so models are int8 and arenas are statically sized
        (\cref{sec:on-device-ai}).
  \item[Connectivity.] The station is provisioned with credentials but never logs in as a user; it
        authenticates to the cloud with a single revocable token (\cref{sec:firmware-station}).
  \item[Hosting.] The dashboard is hosted on a managed Python platform (Streamlit Community Cloud),
        which does not expose server-side cookies, shaping how the session is read
        (\cref{sec:dashboard}).
\end{description}

\Cref{sec:evaluation} revisits these requirements and constraints and reports the extent to which
each is currently met.
